\documentclass[11pt,a4paper]{article}
\usepackage[margin=24mm]{geometry}
\usepackage[T1]{fontenc}
\usepackage{lmodern,amsmath,amssymb,amsthm,booktabs,longtable,array,xcolor,hyperref,microtype,fancyvrb}
\hypersetup{pdftitle={Resultant multiplicities at the fifteen prime Fricke levels},pdfauthor={Research memorandum prepared for David Jordan},colorlinks=true,linkcolor=blue!45!black,urlcolor=blue!45!black}
\newtheorem{theorem}{Theorem}
\newtheorem{lemma}{Lemma}
\setlength{\parindent}{0pt}
\setlength{\parskip}{5pt}
\setlength{\emergencystretch}{2em}
\setlength{\tabcolsep}{4pt}
\newcolumntype{L}[1]{>{\raggedright\arraybackslash}p{#1}}
\title{Resultant multiplicities at the fifteen prime Fricke levels\\[3pt]\large A source audit, geometric proof, and exact computation}
\author{Research memorandum prepared for David Jordan}
\date{28 October 2026}
\begin{document}
\maketitle
\begin{abstract}
The corrected general identity holds in every one of the 1,127 tested pairs, comprising all fourteen odd levels in Nakaya's set and level 2, for every prime $5\le p\le400$, $p\ne N$. A geometric proof for every permitted prime is given below, using norms and divisors on the smooth modular curve; it does not assume that its plane image is smooth. Two qualifications are essential: unramified elliptic factors must be removed before defining $A_N,B_N$, and level 2 needs two different CM exponents. Morton's four original cases are already proved; priority for the extension to the other eleven levels is not established by this literature search.
\end{abstract}

\section{Source transcription and the original four cases}
The following is a transcription of Conjecture~4 in the supplied Nakaya arXiv v1, Section~4.1, p.~18. All equalities are modulo $p$, for primes $p\ge5$, $p\ne N$. Write $s_N=ss_p^{(N*)}(Y)$,
$\delta=(1-(\frac{-3}{p}))/2$, $\epsilon=(1-(\frac{-1}{p}))/2$,
$\nu=(1-(\frac{-2}{p}))/2$, and $\mu_N=(1-(\frac{-N}{p}))/2$.
\begin{align*}
 Y^\epsilon(Y-256)^\nu V_p^{(2*)}(Y)
 &= (Y+144)^{2\delta}(Y-648)^\epsilon s_2^2,\\
 Y^\delta(Y-108)^\delta V_p^{(3*)}(Y)
 &= (Y+192)^{2\delta}(Y^2-576Y-1728)^\epsilon s_3^2,\\
 (Y^2-44Y-16)^{\mu_5}V_p^{(5*)}(Y)
 &= (Y^2+216Y+144)^{2\delta}(Y^2-540Y-6480)^\epsilon s_5^2,\\
 (Y+1)^{\mu_7}(Y-27)^{\mu_7}V_p^{(7*)}(Y)
 &= (Y^2+224Y+448)^{2\delta}\\[-2pt]
 &\quad\cdot(Y^4-528Y^3-9024Y^2-5120Y-1728)^\epsilon s_7^2.
\end{align*}
The page was checked both by text extraction and visual rendering. The code verifies the four displayed correction polynomials explicitly, and verifies the identities at all 302 permitted pairs in these four levels up to 400. There is no discrepancy with Nakaya's formulas.

\subsection{Morton's exact scope}
\begin{center}
\begin{tabular}{llll}
\toprule
Level & Result in supplied Morton v1 & Verified page & Hypothesis\\
\midrule
5 & Theorem 1, equations (6)--(7) & 3 & prime $p\ge7$\\
7 & Theorem 3, equations (13)--(14) & 12 & prime $p\ge5$, $p\ne7$\\
2, 3 & Theorem 5, equations (26)--(27) & 21 & prime $p\ge5$\\
\bottomrule
\end{tabular}
\end{center}
These are exactly all four levels in the original Conjecture~4. For level 5, $p\ge7$ is the same permitted domain as $p\ge5$, $p\ne5$. Morton's exceptional primes in the proofs are checked separately; they are not additional exclusions from the theorems. The introduction, p.~2, explicitly identifies the three theorems as covering all four cases.

\textbf{Pagination limitation.} The uploaded file is arXiv:2105.14532v1, with 27 pages, not the 26-page journal typesetting. The journal citation is \emph{International Journal of Number Theory, 19}(4), 803--828 (2023), DOI 10.1142/S1793042123500409. The theorem numbers and pages in the table were verified in the uploaded preprint. The exact journal page for each theorem, and unchanged numbering in the final typesetting, are \textbf{not verified}; the accessible publisher preview contains only the opening page. No journal page was inferred by adding a fixed offset.

\section{A uniform statement and its proof}
Let $\mathfrak S=\{2,3,5,7,11,13,17,19,23,29,31,41,47,59,71\}$ and use Nakaya's constant terms $104,42,16,9,0,\ldots,0$. Write
\[
 R_N(X,Y)=X^2-a_N(Y)X+b_N(Y),\quad D_N=a_N^2-4b_N,
 \quad Q_{N,c}=R_N(c,Y).
\]
All gcds and factorizations in the following definitions are taken in $\mathbb Q[Y]$, with monic factors. For $c=0,1728$, put
\begin{equation}\label{eq:U}
 U_{N,c}=\prod_{\substack{f\text{ monic irreducible},\ f\mid Q_{N,c}\\
 f\mid a_N-2c}} f^{v_f(Q_{N,c})}.
\end{equation}
Thus $U_{N,c}$ is the full part of $Q_{N,c}$ supported on its gcd with $a_N-2c$, not just that gcd. Define the monic polynomials $A_N,B_N,E_N,H_N$ by
\begin{equation}\label{eq:ABC}
 Q_{N,0}=U_{N,0}A_N^3,\qquad Q_{N,1728}=U_{N,1728}B_N^2,
 \qquad D_N=E_NH_N^2,
\end{equation}
where $E_N$ is squarefree. Equivalently, $E_N$ is the product of the irreducible factors of $D_N$ having odd exponent. These polynomials belong to $\mathbb Z[Y]$: they are monic rational factors of monic integer polynomials, so Gauss's lemma applies.

For an odd prime $p\ge5$, let
\[
 \delta=\frac{1-\left(\frac{-3}{p}\right)}2,\quad
 \epsilon=\frac{1-\left(\frac{-1}{p}\right)}2,\quad
 \mu_N=\frac{1-\left(\frac{-N}{p}\right)}2,\quad \nu=\mu_2.
\]
When $p\ne N$, set
\begin{equation}\label{eq:C}
 C_{N,p}(Y)=
 \begin{cases}
 Y^\epsilon(Y-256)^\nu,&N=2,\\
 E_N(Y)^{\mu_N},&N\text{ odd}.
 \end{cases}
\end{equation}
Bars below mean coefficientwise reduction modulo $p$.

\begin{theorem}[Resultant multiplicities for the prime Fricke levels]\label{thm:main}
Let $N\in\mathfrak S$, and let $p\ge5$ be prime with $p\ne N$. Put
\[
 V_{N,p}(Y)=\operatorname{Res}_X(ss_p(X),\overline{R_N}(X,Y)),
 \qquad S_{N,p}=\operatorname{rad}(V_{N,p}),
\]
where the radical is monic. Then $S_{N,p}=ss_p^{(N*)}$ and
\begin{equation}\label{eq:main}
 \boxed{\quad \overline{C_{N,p}}\,V_{N,p}
       =\overline{A_N}^{\,2\delta}\,
        \overline{B_N}^{\,\epsilon}\,S_{N,p}^{,2}
        \quad\text{in }\mathbb F_p[Y].\quad}
\end{equation}
For every monic irreducible factor $f$ of $V_{N,p}$, its exact exponent is
\begin{equation}\label{eq:exponent}
 v_f(V_{N,p})=2+2\delta\,v_f(\overline{A_N})
                 +\epsilon\,v_f(\overline{B_N})
                 -v_f(\overline{C_{N,p}}).
\end{equation}
In particular, reduction multiplicities in the correction factors must be retained.
If $d_p=\deg ss_p$, then $\deg V_{N,p}=(N+1)d_p$ and its leading coefficient is $1$. There is no unspecified scalar in \eqref{eq:main}.
\end{theorem}

\subsection{The norm lemma, including singular plane points}
\begin{lemma}\label{lem:norm}
Let $k$ be algebraically closed of characteristic different from $2$, let $C$ be a smooth curve with a nonidentity involution $w$, and let $\pi:C\to D=C/\langle w\rangle$. Suppose $Y$ is a local coordinate on $D$, $j$ is regular above the point under consideration, and
\[
 R(X,Y)=(X-j)(X-j\circ w)
       =\operatorname{Norm}_{k(C)/k(D)}(X-j).
\]
For any finite set $J\subset k$, put $s(X)=\prod_{\alpha\in J}(X-\alpha)$. Then
\begin{equation}\label{eq:normorders}
 \operatorname{ord}_{Y=y}\operatorname{Res}_X(s,R)
 =\sum_{P\in\pi^{-1}(y)}\operatorname{ord}_P(s(j))
 =\sum_{\substack{P\in\pi^{-1}(y)\\j(P)\in J}}e_j(P).
\end{equation}
The sums run over distinct points of the normalization, including a fixed point only once.
\end{lemma}
\begin{proof}
The product formula for the resultant gives
$\operatorname{Res}_X(s,R)=\prod_{\alpha\in J}R(\alpha,Y)
=\operatorname{Norm}(s(j))$; the possible signs cancel because the norm has degree $2$.
At an unramified point of $D$ the completed algebra is a product of two power-series rings, and the order of the norm is the sum of the two orders. At a branch point, tame linearization gives parameters $t,u$ with $w(t)=-t$ and $u=t^2$. If $s(j)=t^m v(t)$ with $v$ a unit, its norm is $(-1)^m u^m v(t)v(-t)$, which has order $m$ in $u$. This proves the first equality. Squarefreeness of $s$ gives the second, since the order of $j-j(P)$ is $e_j(P)$. No regularity of the plane curve $R=0$ has been used.
\end{proof}

\subsection{Integral models and the three correction divisors}
\begin{proof}[Proof of Theorem~\ref{thm:main}]
Work over $R=\mathbb Z_{(p)}$. Use the smooth coarse modular curve
$\mathcal C=X_0(N)_R$, its Fricke involution $w$, and its quotient
$\mathcal D=\mathcal C/\langle w\rangle$. Since $p\nmid6N$, the moduli stabilizers and the quotient involution are tame. Jordan's Section~2 and Proposition~5.1 identify $\mathcal D$ with $\mathbb P^1_R$ by $Y=j_N^*$, with its cusp at infinity. The two cusps of $\mathcal C$ are distinct and exchanged by $w$.

On the complement of those cusps, $j$ is regular and $\pi:\mathcal C\to\mathcal D$ is finite locally free of degree two. Its characteristic polynomial is
\[
 \operatorname{Norm}(X-j)=X^2-\operatorname{Tr}(j)X+\operatorname{Norm}(j).
\]
Nakaya's Proposition~4 identifies this polynomial with $R_N$ on the generic fibre. Both sides are integral sections of a torsion-free algebra, so the equality holds over $R$ and after reduction. This is the required specialization statement; it is stronger and safer than assuming that the plane model is smooth.

For completeness, the plane model does have the correct function field in every permitted characteristic. The pole orders of $j$ at the two cusps are $1$ and $N$, whereas those of $j\circ w$ are $N$ and $1$. Their leading coefficients are units, so $j\ne j\circ w$ also after reduction. Hence $k(\mathcal C_k)=k(Y,j)$ and the smooth modular curve is the normalization of $R_N=0$. The map $j$ is separable: away from the elliptic points the level-forgetting morphism is etale. Its degree need not be prime to $p$ for this conclusion.

Let $T_3$ be the reduced divisor of points above $j=0$ with $e_j=3$, and let $T_2$ be the analogous reduced divisor above $j=1728$ with $e_j=2$. At every other finite point $e_j=1$. These descriptions hold on the generic and special fibres. Indeed, over an elliptic curve the choices of a cyclic subgroup of order $N$ form a finite etale scheme; the coarse map is the tame quotient by $\operatorname{Aut}(E)/\{\pm1\}$. A nonfixed subgroup has orbit size $3$ or $2$, respectively; a fixed subgroup has orbit size $1$. This also proves that the ramified and unramified parts of each elliptic fibre specialize with their multiplicities unchanged on $\mathcal C$.

Over characteristic zero, an isogeny with source and target both of invariant $0$, or both of invariant $1728$, is an endomorphism after identifying the curves. The endomorphism ring is commutative, so its kernel is stable under all automorphisms. Conversely, a subgroup stable under all automorphisms gives a quotient with the same exceptional invariant. Thus the support of $U_{N,c}$ in \eqref{eq:U} is precisely the image of the unramified part of $j^{-1}(c)$. Applying Lemma~\ref{lem:norm} to $X$ and $X-1728$ gives
\begin{equation}\label{eq:ramdiv}
 \operatorname{div}_0(A_N)=\pi_*T_3,\qquad
 \operatorname{div}_0(B_N)=\pi_*T_2
\end{equation}
on the generic fibre. These equalities persist on the special fibre. One can see this without dividing polynomials modulo $p$: after an etale base change split the reduced elliptic fibres into sections, and take the products of $Y-Y(P)$ over the appropriate sections. Their monic products are the reductions of the characteristic-zero polynomials in \eqref{eq:ramdiv}. Distinct sections of $\mathcal C$ may have the same image in $\mathcal D$; such repetitions are retained in these products.

Let $F$ be the fixed-point divisor of $w$. Its image on $\mathcal D$ is the branch divisor. In characteristic zero the quadratic extension can be written
\[
 k(\mathcal C)=k(Y)\bigl(\sqrt{D_N(Y)}\bigr)
             =k(Y)\bigl(\sqrt{E_N(Y)}\bigr).
\]
Therefore $\operatorname{div}_0(E_N)=\pi_*F$. There is no branch point at the cusp. Since the involution is tame and nonidentity on every fibre, its fixed scheme is finite etale over $R$: locally at a fixed point the action on a parameter is $t\mapsto-t$. In particular, the branch points remain distinct. This proves that $\overline{E_N}$ is squarefree for every $p\ge5$, $p\ne N$, and that it still cuts out precisely the branch divisor. The factors in $H_N^2$ are not extra branch points.

For odd $N\ge5$, the characteristic-zero Fricke-fixed points have CM by orders of discriminant $-4N$, and also $-N$ when $N\equiv3\pmod4$. Their common CM field is $\mathbb Q(\sqrt{-N})$. Every special-fibre fixed point lifts, by the etaleness just proved. The CM reduction criterion consequently says that all these fixed points are supersingular when $\left(\frac{-N}{p}\right)=-1$, and none are supersingular when that symbol is $1$. Equivalently, one can combine the unchanged total fixed-point count with the supersingular fixed-point count of Beneish and Wen, Proposition~3.16, with $D=N,m=1$; both counts agree in the inert case and the latter is zero in the split case. No ramified case occurs because $p\ne N$.

At $N=3$, the two branch values are $0,108$, with $j$-values $0,54000$, of discriminants $-3,-12$; the same conclusion holds with $\mu_3=\delta$. At $N=2$, the values $0,256$ correspond to discriminants $-4,-8$, respectively, and have different supersingularity tests $\epsilon$ and $\nu$. Thus in all cases the divisor of $\overline{C_{N,p}}$ is exactly the image of the supersingular part of $F$.

Now let $Z$ be the reduced supersingular divisor on $\mathcal C_k$, $k=\overline{\mathbb F}_p$, and let $S$ denote its reduced image on $\mathcal D_k$. Supersingularity is preserved by the degree-$N$ isogeny, so $Z$ is $w$-stable. The orbit count, with a fixed point counted once, is the divisor equality
\begin{equation}\label{eq:orbits}
 \pi_*Z=2S-\pi_*F_{\rm ss}.
\end{equation}
On the other hand, the finite zero divisor of $ss_p(j)$ is
\[
 \operatorname{div}_0(ss_p(j))=Z+2\delta T_3+\epsilon T_2.
\]
Push this equality to $\mathcal D_k$ and apply the norm lemma, \eqref{eq:ramdiv}, and \eqref{eq:orbits}. It gives
\[
 \operatorname{div}_0(V_{N,p})
 =2S+2\delta\operatorname{div}_0(\overline A_N)
       +\epsilon\operatorname{div}_0(\overline B_N)
       -\operatorname{div}_0(\overline C_{N,p}).
\]
The support is precisely the set in Nakaya's definition, so $S$ is the zero divisor of $S_{N,p}=ss_p^{(N*)}$. This proves the claimed equality of divisors, even when correction factors collide. Every polynomial involved is monic. Moreover $R_N(\alpha,Y)$ is monic of degree $N+1$ for each $\alpha$, since $a_N$ and $b_N$ are monic of degrees $N$ and $N+1$. Their product over the $d_p$ supersingular invariants proves the degree, leading coefficient, and scalar-one assertions. Taking the coefficient of an irreducible divisor proves \eqref{eq:exponent}.
\end{proof}

\subsection{What is general, and what is special to these levels}
The norm and divisor argument works for every prime $N\ne p$, with the Fricke quotient in place of the $Y$-line. A single polynomial statement in a Hauptmodul requires genus zero, hence the restriction to $\mathfrak S$ here. This proof uses Jordan's good-reduction coordinate result and tame moduli description, but does not need his nodal genus comparison in Lemma~3.1. That comparison supplies an independent aggregate check, not a local multiplicity calculation.

For odd prime $N$, the numbers of unramified points above the two elliptic invariants are
\[
 u_3=1+\left(\frac{-3}{N}\right),\qquad
 u_2=1+\left(\frac{-1}{N}\right).
\]
At $N=2$ they are $u_3=0,u_2=1$. A formula valid at all prime levels uses the Kronecker symbol $u_2=1+\left(\frac{-4}{N}\right)$ instead. Consequently
\[
 \deg A_N=\frac{N+1-u_3}{3},\qquad
 \deg B_N=\frac{N+1-u_2}{2},\qquad
 \deg E_N=2g(X_0(N))+2.
\]
For odd $N\ge5$, Jordan's fixed-point formula (4.1) gives
\[
 \deg E_N=h_{\rm ord}(-4N)+[N\equiv3\pmod4]h_{\rm ord}(-N)=2d_N.
\]
The same last equality holds for $N=3$ with $d_3=1$. Taking degrees in \eqref{eq:main} therefore recovers Jordan's Corollary~1.3. At $N=2$ it gives
\[
 \deg S_{2,p}=\frac{p-1}{8}+\frac38\left(1-\left(\frac{-1}{p}\right)\right)
                       +\frac14\left(1-\left(\frac{-2}{p}\right)\right).
\]
This is a consequence of the stronger multiplicity formula, not a new claim of priority for the known degree formulas.


\section{Checks on the formulation}
\subsection{Why the two corrections to a naive statement matter}
The polynomial $\operatorname{rad}(D_N)$ usually contains much more than the branch divisor. For example,
\[
 D_2=Y(Y-256)(Y-81)^2.
\]
The factor $Y-81$ is excluded from $E_2$. At this value two distinct points on the normalization can have the same $j$-value. They still contribute separately to the norm. Counting a repeated $X$-root as one point would be wrong.

Also, taking the whole square root of $R_N(1728,Y)$ gives the wrong $B_N$. At level 5,
\[
 R_5(1728,Y)=Y^2(Y^2-540Y-6480)^2,
\]
so $U_{5,1728}=Y^2$ and $B_5=Y^2-540Y-6480$. The omitted square is the image of two unramified elliptic points. The rational gcd definition in \eqref{eq:U} distinguishes it from the ramified square factors.

Finally, the prescription $E_N^{\mu_N}$ cannot be used unchanged at $N=2$. Its branch points have CM fields $\mathbb Q(i)$ and $\mathbb Q(\sqrt{-2})$. For instance at $(N,p)=(2,5)$, $\epsilon=0$, $\nu=1$; multiplying by the additional factor $Y$ would give a false identity. The user's proposed common CM-field description is correct at odd levels, but requires this level-two exception. The factors $Y$ at level 2 and at level 3 lie over unramified exceptional elliptic points, so no extra ramification factor should be inserted for them.

\subsection{Geometric multiplicity rule}
For a supersingular quotient value $y$, let $b(y)$ be 1 if its fibre is Fricke-fixed and 0 otherwise. Let $n_3(y)$ and $n_2(y)$ count the distinct points in that fibre where $e_j=3$ and $e_j=2$. The norm proof gives
\[
 m(y)=2-b(y)+2n_3(y)+n_2(y).
\]
Only supersingular elliptic fibres contribute: $n_3=0$ unless $\delta=1$ and $n_2=0$ unless $\epsilon=1$. A nonfixed fibre has two points, hence can have multiplicity $2,3,4,5,$ or $6$; a fixed fibre has multiplicity $1,2,$ or $3$. In particular, a factor of multiplicity 6 is legitimate, not a failure of the square formula. Formula \eqref{eq:exponent} is its version over $\mathbb F_p$ for arbitrary irreducible factors.

\subsection{Recovery of the four formulas}
\begin{center}\small
\begin{tabular}{ccll}
\toprule
$N$ & $E_N$ & $A_N$ & $B_N$\\
\midrule
2 & $Y(Y-256)$ & $Y+144$ & $Y-648$\\
3 & $Y(Y-108)$ & $Y+192$ & $Y^2-576Y-1728$\\
5 & $Y^2-44Y-16$ & $Y^2+216Y+144$ & $Y^2-540Y-6480$\\
7 & $(Y+1)(Y-27)$ & $Y^2+224Y+448$ & see below\\
\bottomrule
\end{tabular}\end{center}
Here $B_7=Y^4-528Y^3-9024Y^2-5120Y-1728$. Substituting these into \eqref{eq:main}, with $\mu_3=\delta$ and the separate level-two definition \eqref{eq:C}, yields precisely the four identities transcribed above, including the leading scalar 1.

\section{Exact computation}
The final, clean run used PARI/GP 2.15.4 (amd64, GMP 6.3.0) and completed in \textbf{11.624 seconds} of GP wall-clock time on the supplied environment. This excludes installation, literature retrieval, and report generation. No floating-point arithmetic enters any mathematical check. The random seed only fixes finite-field factorization choices.

\begin{center}
\begin{tabular}{rrrrrrr}
\toprule
$N$ & Pairs & Failures & Flagged & $\deg A_N$ & $\deg B_N$ & $\deg E_N$\\
\midrule
2 & 76 & 0 & 3 & 1 & 1 & 2\\
3 & 76 & 0 & 4 & 1 & 2 & 2\\
5 & 75 & 0 & 5 & 2 & 2 & 2\\
7 & 75 & 0 & 8 & 2 & 4 & 2\\
11 & 75 & 0 & 13 & 4 & 6 & 4\\
13 & 75 & 0 & 10 & 4 & 6 & 2\\
17 & 75 & 0 & 14 & 6 & 8 & 4\\
19 & 75 & 0 & 14 & 6 & 10 & 4\\
23 & 75 & 0 & 15 & 8 & 12 & 6\\
29 & 75 & 0 & 18 & 10 & 14 & 6\\
31 & 75 & 0 & 19 & 10 & 16 & 6\\
41 & 75 & 0 & 25 & 14 & 20 & 8\\
47 & 75 & 0 & 27 & 16 & 24 & 10\\
59 & 75 & 0 & 35 & 20 & 30 & 12\\
71 & 75 & 0 & 36 & 24 & 36 & 14\\
\midrule
Total & 1127 & 0 & 246 & & &\\
\bottomrule
\end{tabular}
\end{center}
The fourteen odd levels account for 1,051 pairs; level 2 contributes 76.


\subsection{Reconstruction and independence of the checks}
The code constructs the Hauptmodul as a quotient of a form in the correct Fricke eigenspace by $(\eta(\tau)\eta(N\tau))^k$. It asserts the eigenspace dimension $m+1$, the principal part $q^{-1}$, Nakaya's constant term, and integrality of the retained coefficients. It constructs $a_N,b_N$ by eliminating their pole parts; no stored relation is an input. The positive-coefficient checks implement the two identities of Nakaya's Proposition~4, not a comparison with an external coefficient table (that proposition supplies no such table for all fifteen levels).

The remainder orders range from $O(q^{97})$, $O(q^{96})$ at level 2 to $O(q^{443})$, $O(q^{442})$ at level 71 for the sum and product, respectively. Thus at level 71 all positive coefficients through $q^{442}$ and $q^{441}$ are checked. Principal-part matching certifies the identity once modularity and the pole divisor are known; the extra coefficients are implementation safeguards.

The ordinary supersingular polynomial is computed by a binomial formula and cross-checked, for each of the 76 tested primes, against the radical of the Legendre Hasse-polynomial elimination. Each $V_{N,p}$ is computed by a resultant and fully factored. Its radical is obtained by a derivative/gcd routine with recursive $p$th-root extraction, independently of the predicted correction identity; that radical is also checked against the product of the distinct irreducible factors. This is not an independent moduli construction of $S_{N,p}$: it is precisely Nakaya's definition, as requested.

Every irreducible factor's actual exponent is compared to \eqref{eq:exponent}; the full monic polynomial identity is also checked, which detects any missing support or scalar error. Radical degrees are compared with Jordan's Corollary~1.3 (his level-two formula for $N=2$), and each radical is checked to divide $Y^{p^2}-Y$. In the supplied paper, Corollary~6.3 concerns nonrational quadratic factors, rather than the total degree. This audit uses Corollary~1.3 for the degree, as specified in Task B.3.

\subsection{Collision primes: no exclusions}
The phrase ``a prime dividing $\operatorname{disc}_X R_N$'' needs interpretation: this is a polynomial, monic here, so no prime divides all its coefficients. Also $\operatorname{disc}_Y(D_N)=0$ identically because $D_N$ has square factors. The audit therefore uses the nonzero integer certificates
\[
 \operatorname{disc}(A_N),\ \operatorname{disc}(B_N),\ \operatorname{disc}(E_N),
 \operatorname{Res}(A_N,B_N),\ \operatorname{Res}(A_N,E_N),\ \operatorname{Res}(B_N,E_N),
\]
plus $\operatorname{disc}(\operatorname{rad}_{\mathbb Q}D_N)$. Divisibility of the last integer detects collisions among distinct characteristic-zero discriminant roots, including singular-plane points. These tests concern possible collisions whether or not the affected roots are supersingular.

There are 246 flagged pairs: 239 have a collision in the $A,B,E$ correction data and 7 have only a collision in the larger discriminant-root set. The latter are
\[
 (N,p)=(13,17),(17,13),(19,17),(23,13),(29,13),(41,37),(71,37).
\]
Every one of the 246 identities passes. No permitted prime divides $\operatorname{disc}(E_N)$; exact factorizations of these integers have prime support contained in $\{2,N\}$ for these fifteen levels. Thus two true branch points never coalesce in the conjectured range. A branch point can meet an elliptic correction root, and two ramified points can have the same quotient coordinate.

Two explicit checks illustrate the need to retain reduction multiplicities:
\begin{itemize}
\item $(N,p)=(2,11)$: $V=Y(Y+1)^5$, $S=Y(Y+1)$, $A=B=Y+1$, $C=Y$. Both sides of \eqref{eq:main} equal $Y^2(Y+1)^5$.
\item $(N,p)=(11,5)$: $A=Y^2(Y-1)^2$ modulo 5, $\delta=1$, $\epsilon=0$, $\mu_{11}=0$, and $V=Y^6(Y-1)^6$, $S=Y(Y-1)$. Each of the two factors has predicted and actual exponent $2+2\cdot2=6$.
\end{itemize}
All flagged prime lists appear in Appendix~A; their gcd certificates and the individual factor matrices are in the accompanying files. No smoothness inference is made from the absence of a flag. The proof already accommodates singular plane points and collisions through normalization and divisor specialization.

\subsection{Runs outside the conjecture}
The package includes all 30 pairs with $p=2,3$ and all 13 additional diagonal pairs with $p=N\ge5$, for 43 separate exploratory pairs. For $p=2,3$ the ordinary supersingular polynomial is taken to be $X$. Only the resulting polynomials and radicals are reported: the tame exponents are not asserted there.

For the diagonal odd levels, substituting $\mu_N=1/2$ does not define the original polynomial identity. As a deliberately formal diagnostic, the code checks its squared version
\[
 \overline E_N V_{N,N}^2
 \stackrel{?}{=}\overline A_N^{4\delta}\overline B_N^{2\epsilon}S_{N,N}^4.
\]
It fails at all 13 diagonal levels. For example, at $(N,p)=(5,5)$, $V=(Y+3)^6$, $S=Y+3$, $\overline E_5=(Y+3)^2$ and $\overline A_5=(Y+3)^2$; the two sides have exponents 14 and 12. These are not counterexamples to Theorem~\ref{thm:main}.

\subsection{How to reproduce}
From the extracted \path{nakaya_c4} directory, create \path{results} if necessary and run with \textbf{PARI/GP 2.15.4}:
\begin{Verbatim}[fontsize=\small]
mkdir -p results
gp -fq -s 1000000000 code/verify_c4.gp > results/run.log 2>&1
python3 code/audit_outputs.py
\end{Verbatim}
The driver refuses other GP versions, rather than silently bypassing the reported Hauptmodul dimension issue. This fulfils the pinned-version option; a 2.17.2 fix is not claimed. To change the survey bound, use a small driver file containing
\begin{Verbatim}[fontsize=\small]
P_MAX=1000;
read("code/verify_c4.gp");
\end{Verbatim}
For function-only loading set \path{C4_NO_RUN=1} before \path{read}. The functions \path{nk_haupt}, \path{nk_faber}, \path{nk_relation}, \path{nk_ss}, and \path{nk_rad} adapt Jordan's supplied verification code. The new multiplicity audit, extraction of correction factors, collision certificates, and outside-domain diagnostics are labelled separately in the source. \path{verify_duality.gp} was not present in the uploaded archive.

\section{Prior work and novelty}
This is a bounded literature search, not a proof that no earlier general formula exists. Searches covered the two named papers, visible citing-paper lists, targeted queries for the resultant identity, and the requested Aricheta, Sakai, Tsutsumi, and Kaneko--Zagier literature. Technical claims below rely on the original papers or publisher abstracts, not automated literature summaries.

\begin{longtable}{L{.19\linewidth}L{.45\linewidth}L{.10\linewidth}L{.16\linewidth}}
\toprule Source & What was established or inspected & Overlap & Verification\\\midrule\endhead
Nakaya (2019) & Proposition 4 constructs the two monic coefficient polynomials. Section 4.1, Conjecture 4 states the four identities; Remark 2 identifies elliptic factors. & Partial & Supplied v1, pp. 14, 18, checked.\\\addlinespace
Morton (2023) & Theorems 1, 3, 5 prove the exact four identities at $N=5,7,2,3$, respectively, on their full stated domains. & Full for four; partial overall & Supplied v1, pp. 3, 12, 21; final theorem pagination not verified.\\\addlinespace
Morton (2024) & Tate-form $E_7$ linear/cubic factor counts; quadratic/sextic assertions are explicitly conjectural in the abstract. This is related to Nakaya's Conjecture 5, not an all-level Conjecture 4 proof. & Partial & Original arXiv v5 abstract checked; full theorem audit not verified.\\\addlinespace
Jordan (supplied manuscript) & Section 2 and Proposition 5.1 give the coarse quotient and integral coordinate; Corollaries 1.3 and 6.3 give degrees and quadratic-factor counts. Appendix A reconstructs relations and audits multiple plane points. & Partial & Supplied text and code checked; no claim of publication status.\\\addlinespace
Aricheta (2019) & Theorem 3.1 relates the number of nonrational supersingular points to a genus difference; Remark 1.2 and Appendix A settle the relevant complete-splitting lists. Not the corrected resultant multiplicity formula. & Partial & Original SIGMA paper, pp. 3, 7--8 and relevant sections checked.\\\addlinespace
Hajouji (2023) & Lemmas 2.6--2.7 and Propositions 2.8--2.9 treat automorphism weights, Fricke orbits and $\mathbb F_{p^2}$; Sections 2.4--2.5 use these quadratic relations and resultant radicals and settle Conjecture 2. The fixed-point versus two-point loop distinction is explicit. & Partial & arXiv v3, pp. 6--9 checked. No all-level correction identity found.\\\addlinespace
Beneish \& Wen (2026) & Proposition 3.16 gives the supersingular fixed-point count used above for $D=N,m=1$; Theorem 6.5 gives quotient special fibres and genera. & Partial & Original arXiv v1, p. 15 and Theorem A statement checked. No resultant formula located.\\\addlinespace
Sakai (2011) & Atkin orthogonal polynomials for Fricke levels at most 3 and their relation to supersingular polynomials. & Partial background & Publisher abstract/preview checked; full theorem audit not verified.\\\addlinespace
Sakai (2014) & Atkin polynomials for levels 5 and 7, extremal quasimodular forms, and a conjectural supersingular connection. & Partial background & Publisher abstract checked; full theorem audit not verified.\\\addlinespace
Tsutsumi (2007) & Low-level congruence/triangle-group Atkin polynomials and supersingular elliptic curves. & Partial background & Publisher abstract and metadata checked; full text not verified.\\\addlinespace
Kaneko \& Zagier (1998) & Level-one supersingular polynomials via modular forms, hypergeometric series and Atkin orthogonal polynomials. & Partial background & Author-hosted paper introduction checked; not an all-level Fricke resultant theorem.\\\addlinespace
Nakaya (2023 preprint; 2024 journal article) & Atkin-like orthogonality and generalized Faber expansions at level one. A visible citing work of Morton, but not evidence for the proposed extension. & None to this identity & arXiv introduction checked; full exclusion/ journal pagination not verified.\\
\bottomrule
\end{longtable}

\textbf{Recommendation: a proposition plus a remark.} Place the norm-divisor proposition after Jordan's Proposition~5.1 and move the full computational data into an appendix or supplement. The local norm rule is standard, and the branch/elliptic corrections are a natural consequence once the model and coordinate are set up. The explicit uniform packaging at the eleven additional levels appears useful; an unqualified claim that it is the first such result is not supported by this search. No part of Morton's original four-case theorem should be described as still conjectural.

Suggested wording for Jordan's paper:
\begin{quote}
The same moduli interpretation also determines the multiplicities in Nakaya's resultant. With the unramified elliptic factors removed, the norm-divisor calculation below gives a uniform identity at every prime level in $\mathfrak S$. It recovers the four identities of Nakaya's Conjecture 4, already proved by Morton, and supplies their analogues at the remaining levels. The argument uses standard tame ramification and Fricke fixed-point geometry; we make no priority claim for these geometric ingredients. Exact computations through $p\le400$ are included as checks of the formulas and normalizations.
\end{quote}
The supplied draft has not been modified. The theorem and proof are supplied as a separate editable LaTeX fragment. The proof here is a mathematical derivation from the stated integral modular-curve facts, not a conclusion drawn from finite agreement; no unresolved plane-normalization or collision lemma is being left implicit. The unresolved issues are bibliographic: final journal theorem pagination, inaccessible full texts, and exhaustive priority.

\appendix
\section{All flagged primes in the tested range}
All entries below pass. A prime is listed if any certificate described in Section 4 vanishes modulo that prime.
\begin{longtable}{rp{.86\linewidth}}
\toprule $N$ & Flagged primes $5\le p\le400$, $p\ne N$\\\midrule\endhead
2 & 5, 7, 11\\\addlinespace
3 & 5, 7, 11, 23\\\addlinespace
5 & 7, 11, 19, 59, 71\\\addlinespace
7 & 5, 11, 17, 19, 23, 47, 83, 131\\\addlinespace
11 & 5, 7, 17, 19, 23, 29, 31, 43, 107, 167, 263, 347, 359\\\addlinespace
13 & 5, 7, 11, 17, 23, 43, 47, 107, 251, 311\\\addlinespace
17 & 5, 7, 11, 13, 19, 23, 31, 43, 47, 59, 67, 83, 191, 383\\\addlinespace
19 & 5, 7, 11, 17, 23, 29, 31, 41, 47, 53, 59, 67, 71, 83\\\addlinespace
23 & 5, 7, 11, 13, 17, 19, 31, 41, 43, 47, 53, 67, 79, 83, 107\\\addlinespace
29 & 5, 7, 11, 13, 17, 19, 23, 31, 43, 47, 59, 67, 71, 79, 83, 103, 107, 179\\\addlinespace
31 & 5, 7, 11, 17, 19, 23, 29, 41, 43, 47, 59, 71, 79, 83, 89, 107, 179, 191, 383\\\addlinespace
41 & 5, 7, 11, 17, 19, 23, 29, 31, 37, 43, 47, 59, 67, 71, 79, 83, 103, 107, 127, 131, 139, 151, 163, 239, 383\\\addlinespace
47 & 5, 7, 11, 17, 19, 23, 29, 31, 41, 43, 59, 67, 71, 79, 89, 103, 107, 113, 127, 137, 139, 151, 163, 179, 227, 251, 347\\\addlinespace
59 & 5, 7, 11, 17, 19, 23, 29, 31, 41, 43, 47, 53, 67, 71, 79, 83, 101, 107, 113, 127, 139, 149, 151, 163, 167, 173, 191, 199, 211, 223, 227, 239, 251, 311, 359\\\addlinespace
71 & 5, 7, 11, 17, 19, 23, 29, 31, 37, 41, 43, 47, 53, 59, 67, 79, 83, 89, 101, 103, 113, 127, 131, 137, 139, 149, 163, 167, 197, 199, 211, 223, 227, 239, 271, 283\\\addlinespace
\bottomrule\end{longtable}

\section{Correction-polynomial output and file map}
All fifteen $A_N,B_N,E_N$ are printed in \path{results/run.log}, in
\path{results/correction_polynomials.txt}, and in the exact GP vectors
\path{results/data_N.gp}. The latter have order
\[ [A_N,B_N,E_N,H_N,U_{N,0},U_{N,1728}]. \]
The complete factorizations over $\mathbb Z$ of $R_N(0,Y)$, $R_N(1728,Y)$ and $D_N$ appear in \path{integer_factors_N.txt}. Large coefficients are left in these machine-readable files to avoid transcription errors.

\begin{longtable}{L{.44\linewidth}L{.48\linewidth}}
\toprule File & Contents\\\midrule\endhead
\path{results/pairs.csv} & One row per permitted pair; status, degrees, maximum multiplicity, timing.\\
\path{results/status_matrix.csv} & Requested $N\times p$ PASS table; diagonal marked OUTSIDE.\\
\path{results/factors_N.txt} & Full actual and predicted factor matrices for each prime, including the first mismatch field.\\
\path{results/radicals_N.gp} & All monic supersingular radical polynomials.\\
\path{results/relation_N.gp} & Reconstructed $[a_N,b_N,k,m,v_{\rm sum},v_{\rm prod}]$.\\
\path{results/exceptions.csv} & Collision flags for all 246 affected pairs.\\
\path{results/exceptional_details.txt} & Polynomial gcds explaining each flag. Units in GP gcd output mean no common root.\\
\path{results/certificates.txt} & Exact integer discriminants and pairwise resultants; detects further possible collision primes beyond the survey.\\
\path{results/outside_scope.*} & All 43 exploratory pairs, including full factors and the separate formal diagonal diagnostic.\\
\path{results/audit.json} & Structural checks of the saved output, pair counts, statuses and exception counts.\\
\path{code/verify_c4.gp} & Self-contained reconstruction and multiplicity audit, pinned to GP 2.15.4.\\
\path{report/statement_and_proof.tex} & Editable theorem, definitions, and complete geometric proof.\\
\bottomrule
\end{longtable}

\section*{References}
\addcontentsline{toc}{section}{References}
\footnotesize
\begin{description}
\setlength{\itemsep}{1pt}
\setlength{\parsep}{0pt}
\item Aricheta, V. M. (2019). Supersingular elliptic curves and moonshine. \emph{SIGMA, 15}, Article 007. \url{https://doi.org/10.3842/SIGMA.2019.007}
\item Beneish, L., \& Wen, B. (2026). \emph{Atkin--Lehner quotients of modular curves at primes of bad reduction} [Preprint, v1]. arXiv. \url{https://arxiv.org/abs/2609.39725}
\item Hajouji, N. (2023). \emph{Supersingular isogeny graphs from algebraic modular curves} [Preprint, v3]. arXiv. \url{https://arxiv.org/abs/2303.09096}
\item Jordan, D. (n.d.). \emph{Nakaya's conjectures on the linear factors and degree of Fricke supersingular polynomials} [Unpublished manuscript supplied with this request]. Filename: \path{Nakaya_Conjectures_5_6(3).pdf}. Publication status and date not verified; the associated code archive was also supplied.
\item Kaneko, M., \& Zagier, D. (1998). Supersingular $j$-invariants, hypergeometric series, and Atkin's orthogonal polynomials. In D. A. Buell \& J. T. Teitelbaum (Eds.), \emph{Computational perspectives on number theory} (AMS/IP Studies in Advanced Mathematics, Vol. 7, pp. 97--126). American Mathematical Society. \url{https://www2.math.kyushu-u.ac.jp/~mkaneko/papers/17Atkin.pdf}
\item Morton, P. (2023). Supersingular conjectures for the Fricke group. \emph{International Journal of Number Theory, 19}(4), 803--828. \url{https://doi.org/10.1142/S1793042123500409} [The theorem audit uses the supplied 2021 preprint, \url{https://arxiv.org/abs/2105.14532v1}.]
\item Morton, P. (2024). The Hasse invariant of the Tate normal form $E_7$ and the supersingular polynomial for the Fricke group $\Gamma_0^*(7)$. \emph{The Ramanujan Journal, 63}, 339--385. \url{https://doi.org/10.1007/s11139-023-00764-8} [The abstract's distinction between proved linear/cubic and conjectural quadratic/sextic counts was checked; full paper was not re-audited here. Journal page range follows the supplied Jordan bibliography and was not independently verified.]
\item Nakaya, T. (2019). The number of linear factors of supersingular polynomials and sporadic simple groups. \emph{Journal of Number Theory, 204}, 471--496. \url{https://doi.org/10.1016/j.jnt.2019.04.015} [Formulas checked in supplied arXiv:1809.02363v1.]
\item Nakaya, T. (2023). \emph{On the orthogonality of Atkin-like polynomials and orthogonal polynomial expansion of generalized Faber polynomials} [Preprint]. arXiv. \url{https://arxiv.org/abs/2309.15360} [A 2024 journal version was located; final bibliographic details were not independently audited.]
\item Sakai, Y. (2011). The Atkin orthogonal polynomials for the low-level Fricke groups and their application. \emph{International Journal of Number Theory, 7}(6), 1637--1661. \url{https://doi.org/10.1142/S1793042111004460}
\item Sakai, Y. (2014). The Atkin orthogonal polynomials for the Fricke groups of levels 5 and 7. \emph{International Journal of Number Theory, 10}(8), 2243--2255. \url{https://doi.org/10.1142/S1793042114500766}
\item Tsutsumi, H. (2007). The Atkin orthogonal polynomials for congruence subgroups of low levels. \emph{The Ramanujan Journal, 14}, 223--247. \url{https://doi.org/10.1007/s11139-007-9026-8}
\end{description}
\end{document}
